\documentclass[no-page-number]{nanzan-abstract}

\usepackage{float}
\usepackage{tikz}
\usepackage{xurl}
\usepackage{balance}
\usetikzlibrary{arrows.meta,positioning}

\title{姿勢推定を用いた筋力トレーニング評価システムにおける処理の共通化}
\author{2023TS030 川瀬 翔大}
\professor{宮澤 元}

\begin{document}
\maketitle

\section{背景・先行研究と課題}
カメラ映像から人体の姿勢を推定し，関節位置や角度を利用して運動を評価するシステムが研究されている．これらのシステムでは，映像取得，姿勢推定，ランドマーク取得，特徴量計算，評価・判定，結果提示という複数の処理が連続して行われる．一方，評価対象を複数の運動種目へ拡張すると，種目ごとに評価対象となる関節や基準値などが異なるため，プログラムのどの部分を共通化し，どの部分を種目固有の処理として扱うかが設計上の問題となる．

田邉らは，姿勢推定を用いてサイドレイズやアームカールの運動過程を「最初・途中・終了」の3段階に分類し，運動過程に応じたフィードバックを提示するシステムを提案した\cite{tanabe}．5名を対象に各10回の動作を行った実験では，動作に応じたフィードバックを概ね提示できたことが報告されている．一方，一部の姿勢では判定が難しく，著者らは体格差や動作速度などの影響を考察している\cite{tanabe}．

Mercadal-Baudartらは，単一カメラによる3次元姿勢推定を用いて複数の運動を定量化し，関節角度などの評価指標を用いて運動を評価している\cite{mercadal}．また，PatidarらはMediaPipeとOpenCVを用い，スクワット，アームカール，デッドリフト，サイドレイズなどを対象として，関節角度に基づく評価を行っている\cite{patidar}．これらの研究から，複数種目に対して姿勢情報から評価値を求めることは可能であると考えられる．

しかし，既存研究では，複数種目を扱うシステムを「共通するプログラム」と「種目ごとの設定・処理」に分けたとき，どの処理まで共通化できるかというプログラム構成そのものを本研究の中心的な対象としていない．そこで本研究では，筋力トレーニングを題材として，複数種目に対応する評価システムの処理構成に着目する．

\section{目的}
本研究では，姿勢推定を用いた筋力トレーニング評価システムにおいて，種目ごとに異なる評価項目を設定データとして分離し，姿勢推定から評価・判定までの処理をどこまで共通化できるかを明らかにすることを目的とする．

具体的には，スクワット，アームカール，サイドレイズなどを対象とし，評価対象となる関節や基準値などを変更した場合でも，プログラム本体を変更せずに処理できる構成を実装する．その上で，各処理を「共通処理」「設定変更のみで対応できる処理」「種目ごとに個別実装が必要な処理」に分類し，種目追加時に必要となるプログラム変更を整理する．

\section{研究概要・方法}
本研究で重要とするのは，筋力トレーニングのフォームそのものだけではなく，異なる評価項目を持つ複数種目を一つのシステム構成で扱うためのプログラム設計である．システムを図\ref{fig:flow}のように，姿勢情報を取得する共通部分，種目ごとの設定部分，必要に応じて個別化する処理に分けて考える．

\begin{figure}[H]
  \centering
  \scriptsize
  \begin{tabular}{c}
    \fbox{\parbox{30mm}{\centering Webカメラ\\映像取得}}\\[-1mm]
    $\downarrow$\\[-1mm]
    \fbox{\parbox{30mm}{\centering MediaPipe\\姿勢推定}}\\[-1mm]
    $\downarrow$\\[-1mm]
    \fbox{\parbox{30mm}{\centering ランドマーク\\取得}}\\[-1mm]
    $\downarrow$\\[-1mm]
    \fbox{\parbox{30mm}{\centering 特徴量\\計算}}\\[-1mm]
    $\downarrow$\\[-1mm]
    \fbox{\parbox{30mm}{\centering 設定値との\\比較・判定}}\\[-1mm]
    $\downarrow$\\[-1mm]
    \fbox{\parbox{30mm}{\centering 結果表示・\\フィードバック}}\\[1mm]
    $\uparrow$\\[-1mm]
    \fbox{\parbox{30mm}{\centering 種目別設定\\JSON}}
  \end{tabular}
  \caption{想定するシステム構成（上段：共通処理候補，下段：種目別設定）}
  \label{fig:flow}
\end{figure}

例えば，スクワットでは股関節・膝・足首から膝関節角度を求め，アームカールでは肩・肘・手首から肘関節角度を求める．評価対象となる関節は異なるが，「3点のランドマークから角度を計算する」「計算した値を基準値と比較する」という処理は共通化できる可能性がある．一方，複数の関節を組み合わせた判定や動作の時間的変化など，種目によって処理方法自体が異なる場合は個別処理が必要になる可能性がある．

そこで，種目名，評価対象関節，基準値，許容範囲などをJSONに保持し，プログラム本体から分離する構成を検討する．

\section{実装・確認状況}
現在はPython，OpenCV，MediaPipeを用いて試作を進めている．現時点で確認できているのは，Webカメラによる映像取得，MediaPipeによる姿勢推定，ランドマーク・骨格表示，JSONによる設定値の読み込みである．

一方，関節角度の計算，設定値を用いた比較・判定，JSONによる種目切替，フィードバック表示については，現時点では実装済みとはせず，今後実装・確認する．したがって，現段階では「共通化できた」とは判断せず，共通化を検証するためのシステム構成を具体化している段階である．

\begin{table}[H]
  \caption{中間発表時点での実装・確認状況}
  \label{tbl:status}
  \centering
  \scriptsize
  \begin{tabular}{p{0.64\columnwidth}p{0.23\columnwidth}}
    \hline
    項目 & 状況 \\\hline
    Webカメラによる映像取得 & 確認済み \\
    MediaPipeによる姿勢推定 & 確認済み \\
    ランドマーク・骨格表示 & 確認済み \\
    JSONによる設定値の読み込み & 確認済み \\
    関節角度などの特徴量計算 & 未確認 \\
    設定値との比較・判定 & 未確認 \\
    JSONによる種目切替 & 今後実装 \\
    フィードバック表示 & 未確認 \\
    共通化範囲の評価 & 今後実施 \\\hline
  \end{tabular}
\end{table}

\section{評価計画と考察の観点}
まず，スクワット，アームカール，サイドレイズを候補として，種目ごとに評価対象となる関節と基準値を設定する．同一のプログラムを用いて各種目を処理し，設定データの変更によって評価対象を切り替えられるかを確認する．

評価では，(1) 映像取得，(2) 姿勢推定，(3) ランドマーク取得，(4) 特徴量計算，(5) 基準値との比較，(6) 判定結果の表示，の各処理について，複数種目で共通して利用できるかを調べる．さらに，種目を追加する際にプログラム本体を変更する必要がある箇所を記録する．

結果は，共通処理，設定変更のみで対応可能な処理，個別実装が必要な処理に分類する．これにより，単に「複数種目を動かせたか」ではなく，どの処理を共通化でき，どの処理で種目依存性が発生するのかを明らかにする．なお，設定した角度範囲に入ったことを「正しいフォーム」と同一視せず，あくまで設定した評価基準に対する判定として扱う．

\section{今後の計画}
まず，3点のランドマークから関節角度を計算する処理と，JSONの基準値を用いた比較処理を実装する．次に，スクワット，アームカール，サイドレイズへ同一プログラムを適用し，種目ごとの設定だけで処理を切り替えられるか確認する．

その後，種目を追加した場合に変更が必要となったプログラム箇所を整理し，共通化できる処理と個別実装が必要な処理の境界を明確にする．さらに，角度だけでは扱えない評価項目として，複数関節の組合せや動作の時間的変化を追加した場合にも同じ設計が適用できるかを検討する．最終的には，複数種目を扱う評価システムを設計する際に，どの処理を共通化し，どの情報を設定データとして分離すべきかを整理する．

\section{おわりに}
本研究では，筋力トレーニングを題材として，姿勢推定を利用する評価システムのプログラム構成に着目する．種目ごとに異なる評価項目を設定データとして分離し，姿勢推定から評価・判定までの処理を共通化することで，種目追加時のプログラム変更をどこまで削減できるかを検証する．現在は映像取得，姿勢推定，ランドマーク表示，JSON読み込みまでを確認しており，今後は評価処理を実装して共通化の範囲を実験的に整理する．

\balance
\begin{thebibliography}{9}
\bibitem{tanabe}
田邉英介，椋木雅之，高塚佳代子：運動過程をフィードバックする筋力トレーニング支援システム，2021年度電気・情報関係学会九州支部連合大会講演論文集，p.228，2021．

\bibitem{mercadal}
C. Mercadal-Baudart, C. Liu, G. Farrell, M. Boyne, J. González Escribano, A. Smolic, C. Simms: Exercise quantification from single camera view markerless 3D pose estimation, Heliyon, Vol.10, No.6, e27596, 2024.

\bibitem{patidar}
K. Patidar, M. A. Zaveri: Framework for Gym Posture Estimation Using MediaPipe and OpenCV, 2025 IEEE International Conference on Recent Advances in Computing and Systems (REACS), DOI: 10.1109/REACS67479.2025.11413297, 2025.
\end{thebibliography}

\end{document}
